\documentclass[11pt]{article}
\usepackage[margin=1.1in]{geometry}
\usepackage{times}
\usepackage[T1]{fontenc}
\usepackage{setspace}
\onehalfspacing
\usepackage{parskip}
\pagenumbering{gobble}

\title{Compute as Deterrent}
\author{Flyxion}
\date{September 2026}

\begin{document}
\maketitle

The doctrine begins with a reassuring analogy. Advanced chips are the fissile material of artificial intelligence. Data centers are enrichment facilities. Model weights are warheads. Export controls are nonproliferation treaties. Stability therefore requires the responsible powers to accumulate vastly more compute than everyone else, while insisting the accumulation is defensive.

The analogy fails at the one point that made the original doctrine coherent: verification. Nuclear material is difficult to manufacture, physically conserved, and detectable in transit. Computation is divisible, rentable, concealable, substitutable, and continuously improved by methods designed to do more with less. A weapons inspector can count centrifuges. A compute inspector has to determine whether twelve thousand ordinary machines, three new compression techniques, a leaked checkpoint, and six patient months of optimization amount to a strategic capability.

They do. Just not until Tuesday, and by Tuesday the inspection regime has already been drafted for the machines that existed on Monday.

What survives the transfer is the vocabulary, not the mechanism, and the vocabulary is doing real work. Nobody gets a classified briefing for asking to expand a data center. Plenty of people get one for describing that expansion as strategic deterrence, and the difference in access is the entire reason the framing persists. Compute deterrence deters exactly one adversary: the one considerate enough to build the prohibited system in the shape the regulation imagined. Everyone else routes around it, because there was never a border for it to cross. Call it what the ledger says it is, a capital expenditure with an export license attached, and the case has to be argued on its economics. Call it deterrence, and the case gets to walk in wearing someone else's uniform.

\end{document}
